\subsection{Che cosa dimostrano i risultati}

Il risultato più rilevante della campagna è la differenza tra la GNN statica Stage A e la GNN trajectory-aware Stage D. La prima migliora il funzionale d'interpolazione dei campi gaussiani, ma non trasferisce tale vantaggio alla soluzione PDE del benchmark transiente: nel confronto held-out, gli errori medi risultano superiori a quelli della mesh fissa. La seconda riduce l'errore $L^2$ integrato e finale rispetto alla mesh fissa in tutti i 28 casi considerati e mostra un miglioramento appaiato anche per l'errore in seminorma $H^1$ integrato.

Questa evidenza è coerente con l'ipotesi che la dipendenza temporale della strategia di rilocazione sia rilevante per il problema studiato. Non dimostra tuttavia che la sola architettura trajectory-aware sia responsabile dell'intero guadagno: il risultato appartiene alla pipeline composta da feature dello stato FEM corrente, GNN, filtro di sicurezza, interpolazione temporale dello spostamento e solutore ALE--FEM. Un'ablazione che separi tali componenti è necessaria per attribuire causalmente il beneficio.

\subsection{Il divario tra obiettivo di addestramento e accuratezza della PDE}

I target Stage A sono costruiti minimizzando un errore d'interpolazione di un campo gaussiano, con penalità geometriche e di spostamento. Tale obiettivo non coincide con l'errore della soluzione PDE. La campagna statica mostra direttamente che un miglioramento dell'interpolazione non implica necessariamente una soluzione FEM più accurata. Anche G-Adaptivity ha già affrontato l'ottimizzazione della mesh rispetto all'errore della soluzione agli elementi finiti \citep{rowbottom2025gadaptivity}; pertanto, il contributo del presente studio non può essere formulato come prima integrazione tra GNN e errore FEM. Il problema aperto che emerge dai risultati è più specifico: come apprendere una traiettoria di rilocazione compatibile con uno stato transiente e una formulazione ALE, mantenendo un controllo geometrico deterministico, e come caratterizzare il compromesso tra accuratezza, sicurezza e costo nel solver effettivo.

I risultati Stage D sono incoraggianti, ma la qualità della regressione degli spostamenti target rimane limitata: la validation relative displacement error è pari a 79.43\% e il rapporto tra norma dello spostamento predetto e target è circa 0.62. Inoltre, i target supervisionati sono generati con budget limitato e non sono certificati come ottimi. La discrepanza tra accuratezza di spostamento e accuratezza PDE mostra che il solo errore rispetto al target non è sufficiente a valutare il comportamento del modello, ma non autorizza a concludere che la rete abbia appreso una soluzione ottimale dell'obiettivo fisico.

\subsection{Sicurezza geometrica: efficacia e costo del filtro}

Il filtro deterministico impedisce le inversioni e le degenerazioni osservate dopo l'applicazione nel benchmark, con una piccola discrepanza floating-point rispetto alla soglia stretta in un minimo registrato. La frequenza di intervento Stage D, pari al 75.4\%, resta elevata. Il minimo fattore di contrazione raggiunge valori molto piccoli in ciascuna traiettoria, suggerendo che in alcune finestre la proposta grezza può essere fortemente ridotta. Di conseguenza, l'incremento geometrico effettivamente applicato può essere molto più piccolo di quello predetto dalla rete. La sicurezza non deve essere rimossa per migliorare artificialmente l'accuratezza; è invece necessario migliorare la qualità delle proposte o considerare filtri locali, purché accompagnati da una nuova dimostrazione e da verifiche geometriche.

L'audit non dispone di colonne raw e post-filter indipendenti per tutte le finestre: la duplicazione rilevata in alcuni campi dei CSV impedisce di quantificare completamente il numero di celle non ammissibili prima del filtro sull'intera campagna. Il singolo test di consistenza delle coordinate e i riepiloghi degli interventi costituiscono evidenze utili, ma non sostituiscono una diagnostica completa per finestra. Inoltre, il log della suite di audit riporta un test free-stream/GCL non superato per una differenza di circa $2.52\times10^{-12}$ a fronte di una soglia $10^{-12}$. I residui D-GCL del benchmark e l'esito del test unitario sono quantità distinte: prima della sottomissione il test deve essere riesaminato attraverso il residuo del sistema lineare, senza allentare la tolleranza in modo arbitrario.

\subsection{Costo computazionale e scalabilità}

La pipeline Stage D richiede circa 8.1 volte il tempo della mesh fissa sulla piattaforma utilizzata. L'inferenza GNN e l'aggiornamento delle coordinate rappresentano la maggior parte del costo aggiuntivo; il filtro in sé richiede un tempo trascurabile rispetto a tali componenti. Il risultato non è quindi una strategia di accelerazione nel protocollo attuale, ma una strategia che scambia maggiore tempo di calcolo con una riduzione dell'errore $L^2$.

L'assemblaggio nativo verificato è Galerkin standard e non contiene un termine SUPG; la chiave \texttt{use\_supg} presente in alcuni file YAML non viene consumata dai solver C++ e non rappresenta quindi una stabilizzazione attiva. Non sono stati eseguiti studi sistematici su risoluzioni differenti, hardware diversi o parallelizzazione. Non è dimostrato che il rapporto di costo rimanga costante al crescere del numero di nodi. Anche la generalizzazione rispetto alla risoluzione non può essere dedotta dal fatto che il modello sia rappresentato da un grafo: il numero di nodi, la distribuzione delle feature e la scala degli spostamenti possono cambiare la distribuzione degli input.

\subsection{Limiti sperimentali e validità esterna}

I risultati principali riguardano un dominio quadrato bidimensionale, triangoli $P_1$, una famiglia di soluzioni manufatte gaussiane, coefficienti convettivi e diffusivi fissati e una sola risoluzione iniziale $32\times32$. Il test estrapolativo riguarda soprattutto valori di $\sigma$ esterni all'intervallo di training; non costituisce una dimostrazione di trasferimento a una PDE diversa o a condizioni al contorno non viste. Il numero di casi held-out è 28, sufficiente per un confronto appaiato interno al protocollo ma non per caratterizzare tutte le possibili combinazioni parametriche.

I riferimenti numerici dell'oracolo Stage A sono a budget finito e non sono stazionari secondo il criterio sul gradiente proiettato. I target trajectory-aware Stage D sono generati con un budget ancora più limitato e senza un requisito di convergenza. Le affermazioni del lavoro riguardano quindi le prestazioni dei modelli rispetto alle mesh di riferimento e alle baseline effettivamente calcolate, non la vicinanza a una mesh globalmente ottima.

Infine, la metrica di discrepanza degli incrementi di massa non rappresenta il residuo completo del bilancio discreto della PDE. La D-GCL è verificata separatamente; una verifica fisica del bilancio integrale, comprensiva di sorgente e flussi al contorno, costituisce un'attività distinta. Queste limitazioni devono rimanere esplicite nella versione destinata alla sottomissione.

\subsection{Direzioni di sviluppo}

Le estensioni prioritarie sono: (i) studiare la robustezza rispetto alla risoluzione spaziale e al passo temporale; (ii) valutare una loss di apprendimento più direttamente collegata a un estimatore dell'errore FEM o a un residuo PDE, senza confondere l'obiettivo di training con l'errore di test; (iii) ridurre le contrazioni del filtro migliorando la qualità geometrica delle proposte e documentando raw e post-filter per ogni finestra; (iv) analizzare la scalabilità e il costo dell'inferenza su mesh più grandi; (v) estendere la valutazione a una seconda famiglia di PDE e a geometrie differenti. Tali direzioni sono prospettive, non risultati già dimostrati.
